\section*{Editable literal weights and move labels in original Figures1--2 (transcription)}
\noindent\textbf{Locator convention:} The unchanged original assets retain graph and link geometry. The following is a checked transcription of their printed mathematical labels, not a derived linking matrix or new move proof. Figure1 has a left graph and a right framed-link panel; Figure2 has columns(a),(b),(c), each top and bottom.
\par\medskip
\noindent\textbf{Figure1, left graph:} upper-left endpoint $-2$; central vertex $-1$; lower-left endpoint $-3$; right endpoint $-7$.
\par\noindent\textbf{Figure1, right framed link:} upper-left component $-2$; central component $-1$; lower-left component $-3$; right component $-7$.
\par\medskip
\noindent\textbf{Figure2(a), top, left to right:} $m_1\pm1$, $\pm1$, $m_2\pm1$. Bottom, left to right: $m_1$, $m_2$.
\par\noindent\textbf{Figure2(b), top, left to right:} $m_1\pm1$, $\pm1$. Bottom: $m_1$.
\par\noindent\textbf{Figure2(c), top, left to right:} $m_1$, $0$, $m_2$. Bottom: $m_1+m_2$.
\par\noindent The column markers are $(a)$, $(b)$ and $(c)$; the vertically oriented equivalence sign \rotatebox{-90}{$\simeq$} connects each pair of rows, as shown in the unchanged original asset. Signs and subscript placement above are literal transcriptions; no inferred edge or missing label is added.
